\chapter{Conclusion}
\label{chapter:conclusion}

\begin{introduction}
\begin{flushright}
``Sell me this pen!''
\end{flushright}

\vspace{0.5em}

This chapter concludes the thesis. It recapitulates the research objectives and the four guiding hypotheses, summarizes the principal empirical findings, reflects on their broader implications, and offers a final assessment of the work and its potential extensions.
\end{introduction}


The work described in this thesis set out to investigate whether NLP-derived features, extracted from two complementary textual sources, can enhance the portfolio optimization of sectorial ETFs. The motivation was grounded in a concrete gap in the literature: while the predictive value of sentiment signals has been documented at the individual stock and index level, their systematic integration into sector-level portfolio strategies had remained largely unexplored. Addressing that gap required not only sourcing and processing large volumes of unstructured text but also designing a rigorous experimental framework capable of isolating the contribution of each signal type across multiple optimization paradigms.

Three distinct portfolio optimization frameworks were developed and evaluated under identical experimental conditions. The first was a Deep Reinforcement Learning agent trained via Proximal Policy Optimization (DRL--PPO). The second was a two-stage pipeline combining LSTM-based return prediction with Black-Litterman portfolio construction (LSTM--BL). The third was an end-to-end LSTM model trained to directly maximize the Sharpe ratio (LSTM--SR). Each framework was tested under four feature variants that progressively augmented a technical baseline with sector-level sentiment signals derived from two open sources: the GDELT Project's Global Knowledge Graph, which served as a proxy for institutional financial news, and Reddit discussions sourced from \texttt{r/investing}, \texttt{r/stocks}, and \texttt{r/wallstreetbets}, which captured retail investor sentiment. Sentiment extraction relied on a FinBERT-based pipeline for the news domain and a fine-tuned ModernBERT-based model for Reddit. All frameworks were evaluated out-of-sample over a two-year window spanning July 2023 to June 2025, using cumulative return, Sharpe ratio, volatility, and maximum drawdown as performance metrics.

The study was guided by four formal hypotheses, each of which the empirical results speak to directly. The first conjectured that no single optimization framework would consistently outperform all others across sentiment configurations. This was confirmed: no paradigm dominates universally, and framework choice interacts meaningfully with the type of sentiment introduced, with each approach revealing a distinct risk-return personality. The second and central hypothesis, that NLP-derived sentiment provides incremental value over technical signals alone, received strong support across all three frameworks and both sentiment sources without exception. The third hypothesis predicted that news and social media would contribute in qualitatively distinct ways, with news primarily enhancing risk control and social media primarily driving return generation. The empirical results are consistent with this asymmetry: the two sources do not substitute for one another but rather operate through different channels. The fourth hypothesis, that combining both sources would yield a more balanced profile than either alone, is likewise confirmed: the combined variant consistently sits between the two single-source configurations across all frameworks, suggesting a genuine complementarity rather than redundancy.

Taken together, these findings carry several implications. At the methodological level, they demonstrate that open-source textual data, when processed through domain-adapted language models and integrated into established optimization frameworks, can yield consistent and measurable improvements in risk-adjusted performance. The results are not confined to a single framework or a single textual source, which lends confidence to their generalizability. At the conceptual level, the asymmetry between news and social media signals points toward a more nuanced view of what sentiment data actually captures: institutional news appears to convey risk-relevant information, while retail discourse on social platforms correlates more closely with short-term return dynamics. This distinction has practical consequences for how practitioners might choose to source and weight textual signals depending on their primary portfolio objective.

The methodology developed here, owing to its modular nature and not being inherently tied to the U.S. equity market or to sector ETFs, allows for straightforward extensions. Addressing the mentioned limitations, from richer and more granular textual sources to broader social media coverage, the inclusion of macroeconomic variables, and more adaptive refitting mechanisms, could each deepen the evidence base. The framework could equally be applied to other geographies, asset classes, or equity classification schemes, and its scalability to non-U.S. markets represents a worthwhile avenue for future investigation.

The central question this thesis posed was focused but consequential: can NLP-derived signals from financial news and social media enhance the portfolio optimization of GICS sector ETFs? The answer, grounded in a rigorous out-of-sample empirical evaluation across multiple frameworks and feature configurations, is affirmative. Sentiment features derived from both institutional news and retail social media provide consistent and measurable improvements in risk-adjusted performance relative to technical baselines, the two sources contribute in qualitatively distinct and complementary ways, and their combination tends to produce more balanced and robust portfolios than either source alone. These results contribute a sector-level perspective that has been largely absent from prior work in Financial NLP and offer a practical framework for integrating open-source textual data into systematic portfolio construction, with relevance to both academic research and applied asset management.
